% Paper 6 (GRPO-Registry) cross-pillar synthesis
% NOTE: unlike the empirical pillars, P6 has no external frontier-model
% cross-examination digest; this synthesis ties the resource to the program's
% own cross-pillar evidence and adds no new experimental claims.
\section{Synthesis: Why Each Schema Field Earned Its Place}
\label{sec:p6-synthesis}

The registry's seven run-manifest fields are not a wishlist; each was forced by
a measured failure somewhere in this program. The full reporting standard adds
pass@$k$ as an eighth evaluation requirement; the registry is the container that
makes the run-start lessons reusable:

\begin{itemize}[leftmargin=1.5em, itemsep=0.15em]
  \item \emph{Items 1 and 3 (loss form; backend/precision)} are forced by the
        stack-dependence results: a same-label comparison moved final reward
        $85.6\% \to 5.0\%$, $64\%$ of it from an unreported base-checkpoint
        switch and the rest from other unmatched fields,
        without establishing a causal decomposition. Comparisons that do
        not pin these fields remain confounded.
  \item \emph{Item 4 (per-step ZVF/GU)} is forced by the companion ZVF paper: ZVF
        is ${\sim}0.95$ at both reward extremes in the 368-run audit (the
        interior of that cross-model curve is confounded with configuration),
        so only a trajectory, not an endpoint, is interpretable---and it is cheap to emit, since it needs
        only per-group rewards. That paper's micro-jitter falsification
        (batch ZVF $0.158 \to 0.000$ under $\sim$$10^{-4}$ reward jitter while
        contrast-density measures are invariant) is also why the schema
        records the telemetry \emph{source}: a ZVF of zero means different
        things from different reward pipelines.
  \item \emph{Item 5 (group-size schedule)} is forced by the companion group-size
        paper: expected ZVF-indicator is $p^G + (1-p)^G$, so larger $G$
        mechanically lowers ZVF, and our model$\times G$ sweeps confirm the
        monotone drop (e.g.\ Qwen3-32B: $0.33/0.18/0.08$ at $G{=}4/8/16$).
        Comparing stacks at unrecorded $G$ confounds algorithm with schedule;
        the adaptive-$G$ seed entry shows the field capturing a live
        controller, not just a constant.
  \item \emph{Items 2, 6, 7 (KL; held-out split; decontamination/parser)} are
        forced by the confound audits behind the companion length-bias paper and the benchmark design: silent $\beta$ defaults differ $2\times$
        across framework config dumps. Controlled parser probes and held-out
        evaluation remain requirements, not demonstrated universal remedies.
\end{itemize}

\paragraph{Measured evidence (Section~\ref{sec:p6-measured}).} On the N2
same-stack run (one seed, 40 steps), swapping GRPO for AERO, GIFT, or AREAL
moves last-10 reward by at most $0.020$ (at most $0.011$ over all 40 steps),
while completion length moves more. For the aero arm, which never changed
$G$, the shift is run-to-run variation between identical update rules. We
treat these as one-seed observations, not variant rankings.

Read this way, the registry is the program's exit artifact: the companion papers
established \emph{which} stack facts flip conclusions; the registry is the
minimal data structure in which those facts travel with the stack instead of
with the paper trail.
